\documentclass[11pt,a4paper]{article}
\usepackage[margin=3cm]{geometry}
\usepackage{amsmath,amssymb,amsthm,mathrsfs}
\usepackage[T1]{fontenc}
\usepackage{microtype}
\usepackage[colorlinks=true,linkcolor=blue,citecolor=blue]{hyperref}

\theoremstyle{plain}
\newtheorem{theorem}{Theorem}[section]
\newtheorem{proposition}[theorem]{Proposition}
\newtheorem{lemma}[theorem]{Lemma}
\newtheorem{corollary}[theorem]{Corollary}
\theoremstyle{definition}
\newtheorem{definition}[theorem]{Definition}
\newtheorem{remark}[theorem]{Remark}
\newtheorem{fact}[theorem]{Fact}

\newcommand{\A}{\mathfrak{A}}
\newcommand{\HH}{\mathcal{H}}
\newcommand{\G}{\mathcal{G}}
\newcommand{\R}{\mathbb{R}}
\newcommand{\Z}{\mathbb{Z}}
\newcommand{\Cc}{\mathbb{C}}
\newcommand{\Om}{\Omega}
\newcommand{\ip}[2]{\langle #1,#2\rangle}
\DeclareMathOperator{\spec}{spec}
\DeclareMathOperator{\dist}{dist}
\DeclareMathOperator{\supp}{supp}
\DeclareMathOperator*{\ess}{ess}

\title{Limit points of the finite-volume ground states of $P(\varphi)_2$\\
are ground states of the quasi-local dynamics}
\author{Milestone M0}
\date{}

\begin{document}
\maketitle

\begin{abstract}
Let $H(g)$ be the $P(\varphi)_2$ Hamiltonian with spatial cutoff $g$, let $\Om_g$ be its
(unique) ground state, and let $\omega_g=\ip{\Om_g}{\,\cdot\,\Om_g}$ be the induced state on the
time-zero quasi-local algebra $\A$. Let $\alpha$ be the Glimm--Jaffe automorphism group of $\A$,
which is independent of $g$. We prove that every limit point of the net $(\omega_g)_{g\uparrow 1}$,
in the topology of norm convergence on local algebras, is a \emph{ground state} of the
$C^*$-dynamical system $(\A,\alpha)$.

The proof is an exact identity, not an estimate: it produces no error terms and requires no
hypotheses beyond facts already established by Glimm and Jaffe. Consequently the removal of the
spatial cutoff is reduced, modulo the compactness already available from
\cite[Thm.~2.3]{GJIII}, to the \emph{uniqueness} of the ground state of $(\A,\alpha)$.

We also record two invariance properties obtained at no cost: the set of limit points is stable
under space translations (so uniqueness forces translation invariance, with no space-averaging),
and, for even $P$, every $\omega_g$ is $\Z_2$-invariant.
\end{abstract}

\section{Setting}\label{sec:setting}

\subsection{Free field and local algebras}

Let $\HH_0=\Gamma_s(L^2(\R))$ be the symmetric Fock space over $L^2(\R)$ and
$\mu=(-\partial_x^2+m_0^2)^{1/2}$ with $m_0>0$. For real $f\in H^{-1/2}(\R)$ and real
$h\in H^{1/2}(\R)$ the time-zero fields $\varphi_0(f)$ and $\pi_0(h)$ are essentially
self-adjoint on the finite-particle domain \cite[Prop.~1.2.6]{GJExp}; we denote their closures by
the same symbols and set
\[
  W(f,h):=\exp\bigl(i(\varphi_0(f)+\pi_0(h))\bigr),
\]
a unitary on $\HH_0$. For bounded open $O\subseteq\R$ put
\[
  \A(O):=\{W(f,h)\;:\;f,h\in C_c^\infty(O)\ \text{real}\}''\subseteq B(\HH_0),
  \qquad
  \A:=\overline{\textstyle\bigcup_{O}\A(O)}^{\;\|\cdot\|}.
\]
The net $O\mapsto\A(O)$ is isotonous, so $\A$ is a $C^*$-algebra. For $r>0$ we write
$O_r:=\{x\in\R:\dist(x,O)<r\}$.

Let $\tau_a$ denote the automorphism of $\A$ induced by the space translation
$U_a=\Gamma(e^{ia\,\cdot})$, so $\tau_a(\A(O))=\A(O+a)$.

\subsection{The cutoff Hamiltonian}

Let $P$ be a real polynomial, bounded below, of even degree, and let
\[
  \G:=\bigl\{g\in C_c^\infty(\R)\;:\;0\le g\le 1,\ g\equiv1 \text{ on some nonempty bounded open
  interval}\bigr\}.
\]
For $g\in\G$ set
\[
  H(g):=H_0+\int_\R :\!P(\varphi_0(x))\!:\,g(x)\,dx-E(g),
  \qquad
  E(g):=\inf\spec\Bigl(H_0+\int :\!P(\varphi_0)\!:g\Bigr).
\]

\begin{fact}[Glimm--Jaffe; Rosen]\label{fact:sa}
$H(g)$ is self-adjoint and bounded below, and $H(g)\ge0$ with $\inf\spec H(g)=0$.
\cite{GJI,Rosen70,GJ71}
\end{fact}

\begin{fact}[Glimm--Jaffe]\label{fact:vac}
$0$ is a \emph{simple and isolated} eigenvalue of $H(g)$. Its normalised eigenvector $\Om_g$ may
be chosen strictly positive in the $Q$-space (Schr\"odinger) representation.
\cite[\S2.2--2.3]{GJII}
\end{fact}

We write $\omega_g:=\ip{\Om_g}{\,\cdot\,\Om_g}$, a state on $\A$.

\begin{fact}[Glimm--Jaffe: locality and $g$-independence]\label{fact:loc}
There is a one-parameter group $(\alpha_t)_{t\in\R}$ of $*$-automorphisms of $\A$ such that
\begin{enumerate}
\item[\rm(i)] $\alpha_t(\A(O))\subseteq\A(O_{|t|})$ for every bounded open $O$ and every $t\in\R$;
\item[\rm(ii)] if $A\in\A(O)$ and $g\in\G$ satisfies $g\equiv1$ on $O_{|t|}$, then
      $\alpha_t(A)=e^{itH(g)}Ae^{-itH(g)}$.
\end{enumerate}
In particular $\alpha$ does not depend on $g$. \cite[\S III]{GJI}, \cite[\S3.5--3.6]{GJII},
\cite[\S4.1]{GJExp}
\end{fact}

Property (i) is an \emph{exact} light cone: there is no tail, in contrast with the Lieb--Robinson
bounds available for lattice systems. This is used decisively in \S\ref{sec:spectral}.

\subsection{The directed set}

Order $\G$ by
\[
  g\preccurlyeq g' \quad:\Longleftrightarrow\quad \{g=1\}^\circ\subseteq\{g'=1\}^\circ ,
\]
and for bounded open $O$ put
\[
  d(g,O):=\dist\bigl(O,\ \R\setminus\{g=1\}^\circ\bigr)\in(0,\infty].
\]

\begin{lemma}\label{lem:directed}
$(\G,\preccurlyeq)$ is directed; it is stable under translations, and for each $a\in\R$ the map
$g\mapsto\tau_ag$ is an order isomorphism of $\G$ with cofinal range. Moreover, for every bounded
open $O$ and every $R>0$ there is $g_0\in\G$ with $d(g,O)>R$ for all $g\succcurlyeq g_0$.
\end{lemma}

\begin{proof}
Given $g_1,g_2\in\G$, the sets $\{g_i=1\}^\circ$ are bounded and open; choose a bounded open
interval $I$ containing their closures and $g_3\in C_c^\infty(\R)$ with $0\le g_3\le1$ and
$g_3\equiv1$ on $I$. Then $g_3\in\G$ and $g_3\succcurlyeq g_1,g_2$. Translation stability and the
order-isomorphism property are immediate from $\{\tau_ag=1\}^\circ=\{g=1\}^\circ+a$; cofinality
follows from directedness. For the last claim choose $g_0\equiv1$ on a bounded open interval
containing $\overline{O_R}$; then $g\succcurlyeq g_0$ gives
$\{g=1\}^\circ\supseteq O_R$, hence $d(g,O)\ge R$; enlarging $g_0$ slightly gives the strict
inequality.
\end{proof}

\section{Ground states of $(\A,\alpha)$}

Because $\A$ contains Weyl unitaries, $t\mapsto\alpha_t(A)$ is in general \emph{not}
norm-continuous: two Weyl unitaries $W(f,h)$, $W(f',h')$ with $(f,h)\neq(f',h')$ are at distance
$2$ in the CCR $C^*$-algebra. Hence $\alpha$ has no densely defined norm-generator and the
familiar characterisation ``$-i\,\omega(A^*\delta(A))\ge0$'' is unavailable. We therefore use the
spectral form throughout. (See Remark~\ref{rem:nogenerator}.)

\begin{definition}\label{def:gs}
A state $\omega$ on $\A$ is a \emph{ground state} of $(\A,\alpha)$ if
\begin{enumerate}
\item[\rm(a)] $\omega\circ\alpha_t=\omega$ for all $t\in\R$; and
\item[\rm(b)] $t\mapsto\omega(A\,\alpha_t(B))$ is continuous for all $A,B\in\A$, and for every
      $f\in L^1(\R)$ whose inverse Fourier transform
      \[
        \check f(E):=\int_\R f(t)\,e^{itE}\,dt
      \]
      vanishes on $[0,\infty)$, one has $\displaystyle\int_\R f(t)\,\omega\bigl(A\,\alpha_t(B)\bigr)\,dt=0$.
\end{enumerate}
\end{definition}

The continuity requirement in (b) is what makes the integral meaningful; it is verified
independently in Theorem~\ref{thm:main}, so it is not an extra assumption there.

\begin{lemma}\label{lem:gsGNS}
Let $\omega$ satisfy \textup{(a)}. Then \textup{(b)} holds if and only if, in the GNS triple
$(\HH_\omega,\pi_\omega,\Om_\omega)$, the unitary group $U$ implementing $\alpha$ with
$U(t)\Om_\omega=\Om_\omega$ has self-adjoint generator $H_\omega\ge0$. In particular
Definition~\ref{def:gs} agrees with the usual notion \cite[Def.~5.3.18]{BR2}.
\end{lemma}

\begin{proof}
By (a) the group $U(t)\pi_\omega(A)\Om_\omega:=\pi_\omega(\alpha_t(A))\Om_\omega$ is well defined,
unitary and weakly, hence strongly, continuous; write $U(t)=e^{itH_\omega}$. For $A,B\in\A$ set
$u=\pi_\omega(A)^*\Om_\omega$, $v=\pi_\omega(B)\Om_\omega$ and let $\rho:=\ip{u}{P_{H_\omega}(\cdot)v}$,
a complex measure on $\R$ of finite total variation $|\rho|(\R)\le\|u\|\|v\|$. Using
$U(t)^*\Om_\omega=\Om_\omega$,
\[
  \omega\bigl(A\alpha_t(B)\bigr)=\ip{u}{U(t)v}=\int_\R e^{itE}\,d\rho(E).
\]
Fubini (legitimate since $f\in L^1$ and $|\rho|(\R)<\infty$) gives
$\int f(t)\,\omega(A\alpha_t(B))\,dt=\int_\R\check f(E)\,d\rho(E)$.
If $H_\omega\ge0$ then $\rho$ is supported in $[0,\infty)$, where $\check f$ vanishes, so the
integral is $0$. Conversely, if the integral vanishes for all such $f$, then, since
$\{\check f: f\in L^1,\ \check f|_{[0,\infty)}=0\}\supseteq C_c^\infty((-\infty,0))$, we get
$\rho|_{(-\infty,0)}=0$ for all $A,B$; as vectors of the form $\pi_\omega(B)\Om_\omega$ are dense,
$P_{H_\omega}((-\infty,0))=0$.
\end{proof}

\section{Two invariances, for free}\label{sec:inv}

\begin{lemma}[local $\alpha$-invariance]\label{lem:localinv}
Let $O$ be bounded open, $A\in\A(O)$, $g\in\G$ and $|t|<d(g,O)$. Then
$\omega_g(\alpha_t(A))=\omega_g(A)$.
\end{lemma}

\begin{proof}
Since $|t|<d(g,O)$ we have $O_{|t|}\subseteq\{g=1\}^\circ$, so $g\equiv1$ on $O_{|t|}$ and
Fact~\ref{fact:loc}(ii) gives $\alpha_t(A)=e^{itH(g)}Ae^{-itH(g)}$. As $H(g)\Om_g=0$ we have
$e^{-itH(g)}\Om_g=\Om_g$, whence
$\omega_g(\alpha_t(A))=\ip{\Om_g}{e^{itH(g)}Ae^{-itH(g)}\Om_g}=\ip{\Om_g}{A\Om_g}$.
\end{proof}

\begin{lemma}[translation covariance of the net]\label{lem:transl}
For every $a\in\R$ and $g\in\G$ one has $\omega_{\tau_ag}=\omega_g\circ\tau_{-a}$.
\end{lemma}

\begin{proof}
$U_a$ implements $\tau_a$ and commutes with $H_0$, while
$U_a\bigl(\int:\!P(\varphi_0)\!:g\bigr)U_a^*=\int:\!P(\varphi_0)\!:\,\tau_ag$. Hence
$H(\tau_ag)=U_aH(g)U_a^*$ and $E(\tau_ag)=E(g)$. By Fact~\ref{fact:vac} the ground state of
$H(\tau_ag)$ is $U_a\Om_g$ up to a phase, which does not affect the induced state. Therefore
$\omega_{\tau_ag}(A)=\ip{U_a\Om_g}{AU_a\Om_g}=\ip{\Om_g}{U_a^*AU_a\Om_g}
=\omega_g(\tau_{-a}(A))$.
\end{proof}

\begin{lemma}[$\Z_2$ invariance]\label{lem:z2}
Assume $P$ is even and let $\theta$ be the automorphism of $\A$ with $\theta(W(f,h))=W(-f,-h)$.
Then $\omega_g\circ\theta=\omega_g$ for every $g\in\G$.
\end{lemma}

\begin{proof}
$\theta$ is implemented by $\Theta=\Gamma(-1)$, which is unitary, self-adjoint, commutes with
$H_0$, and satisfies $\Theta:\!P(\varphi_0(x))\!:\Theta^*=\;:\!P(\varphi_0(x))\!:$ because $P$ is
even. Hence $\Theta H(g)\Theta^*=H(g)$, so $\Theta\Om_g$ is a normalised ground state and, by
simplicity, $\Theta\Om_g=c\,\Om_g$ with $c\in\{\pm1\}$ (using $\Theta^2=1$). In the $Q$-space
representation $\Theta$ acts by $(\Theta\psi)(q)=\psi(-q)$, so $\Theta\Om_g>0$ pointwise by
Fact~\ref{fact:vac}; hence $c=+1$ and
$\omega_g(\theta(A))=\ip{\Theta\Om_g}{A\Theta\Om_g}=\omega_g(A)$.
\end{proof}

\begin{remark}\label{rem:noaveraging}
Lemma~\ref{lem:transl} replaces the space-translation averaging used in
\cite[eq.~(2.5)]{GJIII}: the set of limit points of $(\omega_g)_{g\in\G}$ is stable under
translations by Lemmas~\ref{lem:directed} and \ref{lem:transl}, so \emph{uniqueness of the limit
already forces translation invariance}. This matters beyond aesthetics: averaging is precisely
what produces a mixed limit in a multiple-phase regime, whereas Lemma~\ref{lem:z2} shows the
unaveraged net is $\Z_2$-invariant by itself.
\end{remark}

\section{The exact spectral identity}\label{sec:spectral}

\begin{proposition}\label{prop:exact}
Let $O$ be bounded open, $A,B\in\A(O)$, $g\in\G$, and set $d:=d(g,O)$, $M:=\|A\|\,\|B\|$. Define
the complex measure
\[
  \rho_g(\cdot):=\ip{A^*\Om_g}{\,P_{H(g)}(\cdot)\,B\Om_g}
\]
on $\R$, and $G_g(z):=\int e^{izE}\,d\rho_g(E)$ for $\operatorname{Im}z\ge0$. Then
\begin{enumerate}
\item[\rm(i)] $\supp\rho_g\subseteq[0,\infty)$ and $|\rho_g|(\R)\le M$;
\item[\rm(ii)] $G_g$ is holomorphic on $\{\operatorname{Im}z>0\}$, continuous and bounded by $M$
      on $\{\operatorname{Im}z\ge0\}$;
\item[\rm(iii)] $\omega_g\bigl(A\,\alpha_t(B)\bigr)=G_g(t)$ for every $|t|<d$;
\item[\rm(iv)] $\displaystyle\int_\R f(t)\,G_g(t)\,dt=0$ for every $f\in L^1(\R)$ with
      $\check f|_{[0,\infty)}=0$.
\end{enumerate}
\end{proposition}

\begin{proof}
(i) $\supp\rho_g\subseteq\spec H(g)\subseteq[0,\infty)$ by Fact~\ref{fact:sa}. For a complex
spectral measure, $|\rho_g|(\R)\le\|A^*\Om_g\|\,\|B\Om_g\|\le M$.

(ii) For $\operatorname{Im}z\ge0$ and $E\ge0$ we have $|e^{izE}|=e^{-E\operatorname{Im}z}\le1$;
holomorphy and continuity follow from dominated convergence together with (i).

(iii) Let $|t|<d$. Then $g\equiv1$ on $O_{|t|}$, so Fact~\ref{fact:loc}(ii) applies to $B$ and
\[
  \omega_g\bigl(A\alpha_t(B)\bigr)
  =\ip{\Om_g}{A\,e^{itH(g)}Be^{-itH(g)}\Om_g}
  =\ip{\Om_g}{A\,e^{itH(g)}B\,\Om_g}
  =\ip{A^*\Om_g}{e^{itH(g)}B\Om_g}
  =G_g(t),
\]
where the second equality uses $e^{-itH(g)}\Om_g=\Om_g$ and the last is the spectral theorem.

(iv) By (i) and $f\in L^1$, Fubini applies:
$\int f(t)G_g(t)\,dt=\int_{[0,\infty)}\check f(E)\,d\rho_g(E)=0$, since $\check f$ vanishes on
$[0,\infty)$.
\end{proof}

The cancellation $e^{-itH(g)}\Om_g=\Om_g$ in (iii) is the whole point: it is exact, and it is
available because $\Om_g$ is the ground state of the \emph{same} operator that implements
$\alpha_t$ locally. No approximation, and no error term, enters.

\begin{theorem}\label{thm:main}
Let $(g_n)_{n\in\mathbb N}\subseteq\G$ be a sequence such that $d(g_n,O)\to\infty$ for every
bounded open $O$, and suppose that for every bounded open $O$ the restrictions
$\omega_{g_n}|_{\A(O)}$ converge in \emph{norm} in $\A(O)^*$. Let $\omega$ be the resulting state
on $\A$. Then $\omega$ is a ground state of $(\A,\alpha)$ in the sense of
Definition~\ref{def:gs}.
\end{theorem}

\begin{proof}
\emph{(a)} Fix $A\in\A(O)$ and $t\in\R$. For $n$ large, $d(g_n,O)>|t|$, so
Lemma~\ref{lem:localinv} gives $\omega_{g_n}(\alpha_t(A))=\omega_{g_n}(A)$. Both
$\alpha_t(A)\in\A(O_{|t|})$ and $A\in\A(O_{|t|})$, so letting $n\to\infty$ yields
$\omega(\alpha_t(A))=\omega(A)$. As $\bigcup_O\A(O)$ is norm-dense in $\A$ and $\|\alpha_t\|=1$,
(a) follows.

\emph{(b)} Fix $A,B\in\A(O)$ and put $M:=\|A\|\|B\|$.

\emph{Continuity.} Let $T>0$ and let $n$ be large enough that $d(g_n,O)>T$. For $|t|\le T$ we
have $A\in\A(O)\subseteq\A(O_T)$ and, by Fact~\ref{fact:loc}(i),
$\alpha_t(B)\in\A(O_{|t|})\subseteq\A(O_T)$; hence $A\alpha_t(B)\in\A(O_T)$ with
$\|A\alpha_t(B)\|\le M$, and therefore
\begin{equation}\label{eq:unif}
  \sup_{|t|\le T}\bigl|\omega\bigl(A\alpha_t(B)\bigr)-\omega_{g_n}\bigl(A\alpha_t(B)\bigr)\bigr|
  \ \le\ M\,\bigl\|(\omega-\omega_{g_n})|_{\A(O_T)}\bigr\|\ \xrightarrow[n\to\infty]{}\ 0 .
\end{equation}
By Proposition~\ref{prop:exact}(ii)--(iii), $t\mapsto\omega_{g_n}(A\alpha_t(B))=G_{g_n}(t)$ is
continuous on $[-T,T]$. Thus $t\mapsto\omega(A\alpha_t(B))$ is a uniform limit of continuous
functions on $[-T,T]$, hence continuous there; as $T$ was arbitrary it is continuous on $\R$, and
bounded by $M$. In particular the integral below is well defined.

\emph{Vanishing.} Let $f\in L^1(\R)$ with $\check f|_{[0,\infty)}=0$. For $T>0$ set
$\varepsilon(T):=\int_{|t|>T}|f(t)|\,dt$, so $\varepsilon(T)\to0$ as $T\to\infty$. Fix $T>0$ and
decompose
\[
  \int_\R f\,\omega(A\alpha_\cdot(B))
  =\underbrace{\int_{|t|\le T} f\,\bigl[\omega-\omega_{g_n}\bigr](A\alpha_t(B))\,dt}_{(\mathrm I)}
  +\underbrace{\int_{|t|\le T} f\,\omega_{g_n}(A\alpha_t(B))\,dt}_{(\mathrm{II})}
  +\underbrace{\int_{|t|> T} f\,\omega(A\alpha_t(B))\,dt}_{(\mathrm{III})}.
\]

By \eqref{eq:unif},
$|(\mathrm I)|\le\|f\|_{L^1}\,M\,\|(\omega-\omega_{g_n})|_{\A(O_T)}\|$,
and $|(\mathrm{III})|\le M\varepsilon(T)$ since $|\omega(A\alpha_t(B))|\le M$.

For $(\mathrm{II})$: choose $n$ so large that $d(g_n,O)>T$. Then
Proposition~\ref{prop:exact}(iii) applies for all $|t|\le T$, so, using
Proposition~\ref{prop:exact}(iv) and then (ii),
\[
  (\mathrm{II})=\int_{|t|\le T} f\,G_{g_n}
  =\underbrace{\int_\R f\,G_{g_n}}_{=\,0}-\int_{|t|>T} f\,G_{g_n},
  \qquad |(\mathrm{II})|\le M\varepsilon(T).
\]

Letting $n\to\infty$ with $T$ fixed kills $(\mathrm I)$ by hypothesis, and we obtain
\[
  \Bigl|\int_\R f(t)\,\omega\bigl(A\alpha_t(B)\bigr)\,dt\Bigr|\ \le\ 2M\varepsilon(T)
  \qquad\text{for every }T>0 .
\]
Letting $T\to\infty$ gives $0$. Finally, for general $A,B\in\A$ choose local
$A_k\to A$, $B_k\to B$ in norm; the map $(A,B)\mapsto\int f\,\omega(A\alpha_\cdot(B))$ is
jointly norm-continuous with bound $\|f\|_{L^1}\|A\|\|B\|$, so (b) extends to $\A$.
\end{proof}

\begin{corollary}\label{cor:exists}
Assume the local number bound
\[
  \sup_{g\in\G}\ \omega_g\bigl(N_{\tau,B}\bigr)\le c(|B|)<\infty \tag{H2}
\]
so that \cite[Thm.~2.3]{GJIII} applies. Then sequences as in Theorem~\ref{thm:main} exist: any
sequence $(g_n)\subseteq\G$ with $d(g_n,O)\to\infty$ for all $O$ admits a subsequence along which
$\omega_{g_n}|_{\A(O)}$ converges in norm for every bounded open $O$, and every resulting limit is
a locally normal ground state of $(\A,\alpha)$.
\end{corollary}

\begin{proof}
Fix an exhausting sequence $O_1\subseteq O_2\subseteq\cdots$ of bounded open sets with
$\bigcup_kO_k=\R$. By \cite[Thm.~2.3]{GJIII} and (H2), $\{\omega_{g_n}|_{\A(O_k)}\}_n$ is
contained in a norm-compact subset of $\A(O_k)^*$ for each $k$; norm-compact metric spaces are
sequentially compact, so a diagonal argument produces a subsequence converging in norm on every
$\A(O_k)$, hence on every $\A(O)$ by isotony. Limits of norm-convergent sequences of normal
states are normal \cite[Thm.~2.3]{GJIII}. Theorem~\ref{thm:main} applies.
\end{proof}

\begin{corollary}[reduction of the target]\label{cor:reduction}
Assume \textup{(H2)}. If $(\A,\alpha)$ has a \emph{unique} locally normal ground state
$\omega_\infty$ which is $\Z_2$- and translation-invariant, then
\[
  \omega_g\longrightarrow\omega_\infty\qquad\text{in norm on every }\A(O),
\]
along the full net $\G$.
\end{corollary}

\begin{proof}
Suppose not; then there are $O$, $\varepsilon>0$ and a sequence $g_n$, cofinal in $\G$, with
$\|(\omega_{g_n}-\omega_\infty)|_{\A(O)}\|\ge\varepsilon$. By Corollary~\ref{cor:exists} pass to a
subsequence converging in norm on all local algebras to a locally normal ground state
$\omega_2$; by Lemmas~\ref{lem:transl}, \ref{lem:z2} and Remark~\ref{rem:noaveraging}, $\omega_2$
is $\Z_2$-invariant, and by hypothesis $\omega_2=\omega_\infty$. This contradicts
$\|(\omega_2-\omega_\infty)|_{\A(O)}\|\ge\varepsilon$.
\end{proof}

\section{Remarks}

\begin{remark}[why the generator form is unusable]\label{rem:nogenerator}
The characterisation of ground states by $-i\,\omega(A^*\delta(A))\ge0$ requires
$t\mapsto\alpha_t(A)$ to be norm-differentiable on a dense $*$-subalgebra. For a Weyl algebra
this fails badly: $\|W(f,h)-W(f',h')\|=2$ whenever $(f,h)\neq(f',h')$, so $\alpha$ is not even
norm-continuous and no densely defined norm-generator exists. Definition~\ref{def:gs} and
Proposition~\ref{prop:exact} avoid generators entirely, working only with bounded operators and
the unitary group $e^{itH(g)}$. As a byproduct, the domain hypothesis
``$A\Om_g\in D(H(g)^{1/2})$'', which a generator-based argument would need, is not required.
\end{remark}

\begin{remark}[M0 needs P2: the two pillars are not independent]\label{rem:dependency}
Theorem~\ref{thm:main} assumes \emph{norm} convergence on local algebras, not merely weak-$*$
convergence of $(\omega_{g_n})$. This is not a technical convenience. In the proof, term
$(\mathrm I)$ must be shown to vanish, and weak-$*$ convergence gives only pointwise convergence
in $t$ of the integrand, with no dominated convergence theorem available for nets. Norm
convergence on $\A(O_T)$ supplies exactly the uniformity in $t\in[-T,T]$ that is needed. Hence
the compactness of \cite[Thm.~2.3]{GJIII} enters the argument \emph{twice}: once to produce a
limit, once to pass the spectral condition to it.
\end{remark}

\begin{remark}[what is \emph{not} proved here]\label{rem:notproved}
Theorem~\ref{thm:main} does not assert that a limit exists along the full net; it identifies the
nature of every limit point. Nor does it produce a rate. The remaining inputs are:
\textup{(H2)} without space-averaging (Milestone M1); uniqueness of the locally normal
$\Z_2$- and translation-invariant ground state of $(\A,\alpha)$ (Milestone M3, or the
uniform-gap substitute); and, for rates, a uniform spectral gap together with a quasi-adiabatic
argument exploiting the exact light cone of Fact~\ref{fact:loc}(i) (Milestone M4).
\end{remark}

\begin{remark}[what M0 buys]
Rosen's proof that $\omega_g\to\omega_\infty$ in the quadratic model \cite[\S5]{Rosen72} has two
steps: identify the limit, then upgrade the topology using \cite[Thm.~4.1]{GJIII}. In the
quadratic model the first step is free because $\mu_g\to\mu_1$ in the one-particle space. For
$P(\varphi)_2$ it is not, and this is exactly what Guerra--Rosen--Simon list as open
\cite[Problems 1--2]{GRS75}. Theorem~\ref{thm:main} supplies that first step in a form that is
exact and independent of any Euclidean input, and thereby reduces the removal of the spatial
cutoff to a \emph{uniqueness} question, where the post-1989 literature has delivered.
\end{remark}

\begin{thebibliography}{99}

\bibitem{BR2} O.~Bratteli and D.~W.~Robinson,
\emph{Operator Algebras and Quantum Statistical Mechanics 2}, 2nd ed., Springer, 1997.

\bibitem{GJI} J.~Glimm and A.~Jaffe,
\emph{A $\lambda\phi^4$ quantum field theory without cutoffs. I},
Phys.\ Rev.\ \textbf{176} (1968) 1945--1951.

\bibitem{GJII} J.~Glimm and A.~Jaffe,
\emph{The $\lambda(\phi^4)_2$ quantum field theory without cutoffs. II. The field operators and
the approximate vacuum}, Ann.\ of Math.\ \textbf{91} (1970) 362--401.

\bibitem{GJIII} J.~Glimm and A.~Jaffe,
\emph{The $\lambda(\phi^4)_2$ quantum field theory without cutoffs. III. The physical vacuum},
Acta Math.\ \textbf{125} (1970) 203--261.

\bibitem{GJ71} J.~Glimm and A.~Jaffe,
\emph{Positivity and self adjointness of the $P(\phi)_2$ Hamiltonian},
Comm.\ Math.\ Phys.\ \textbf{22} (1971) 253--258.

\bibitem{GJExp} J.~Glimm and A.~Jaffe,
\emph{Quantum Field Theory and Statistical Mechanics: Expositions}, Birkh\"auser, 1985.

\bibitem{GRS75} F.~Guerra, L.~Rosen and B.~Simon,
\emph{The $P(\phi)_2$ Euclidean quantum field theory as classical statistical mechanics},
Ann.\ of Math.\ \textbf{101} (1975) 111--189 and 191--259.

\bibitem{Rosen70} L.~Rosen,
\emph{A $\lambda\phi^{2n}$ field theory without cutoffs},
Comm.\ Math.\ Phys.\ \textbf{16} (1970) 157--183.

\bibitem{Rosen72} L.~Rosen,
\emph{Renormalization of the Hilbert space in the mass shift model},
J.\ Math.\ Phys.\ \textbf{13} (1972) 918--927.

\end{thebibliography}

\end{document}
